\section{模型评价}
\subsection{模型优点}
本文模型的首要优点是变量链完整。问题一的每个重叠率均可追溯到位置、水深和两侧覆盖半宽；问题二明确分离航向水深变化与横航向覆盖变化；问题三将边界、重叠和停止条件化为显式递推；问题四则把地形、条带、路线和评估拆成四个接口清晰的模块。这样的结构便于复算，也避免在不同小问中重复或冲突地定义同一物理量。

第二，算法描述包含搜索区间、迭代精度和异常处理。问题三不依赖黑箱求解器，问题四的二分次数、样点数量、曲线细分和敏感性网格均明确给出。原候选 C 因硬拐角失败后，并未直接用离散折线掩盖问题，而是建立显式 Bézier 几何并重新计算全部指标。

第三，结果分析直接对应题目要求。问题一指出具体漏测位置；问题二解释矩阵的方向性；问题三给出条数、总长度和路线族内最少性；问题四给出总长度、漏测比例和过度重叠长度，同时保留保守回退。模型检验中保留 $70\times70$ 网格的不利结果，使结论强度与证据相匹配。

\subsection{模型局限}
问题四的覆盖面积由有限网格中心采样估计，无法替代连续几何上的严格覆盖证明；虽然三组网格能够识别一定的分辨率风险，但仍可能遗漏更小尺度缺口。双线性插值在网格单元内只保留一次变化，会平滑原始采样之外的局部起伏。候选路线族也并非所有可能曲线的全集，因此 C-05 只能称为已评估候选中的综合推荐，而不能称为全局最短。

此外，模型未包括声速剖面、姿态补偿、障碍物、天气、调头连接和具体船型动力学。曲率有限只是路线几何的必要条件，不是实际航行可执行性的充分条件。$L_{20+}$ 采用“后一条线相对前一条线”的有序计数政策，适合回答题目规定指标，但与按物理面积计算的重叠量并不相同，使用结果时必须保留该定义。

\subsection{改进方向}
第一，可对条带边界使用自适应网格加密或区间算术，直接计算覆盖边界的上下界，从而把采样判断升级为更严格的连续覆盖证据。第二，可在 C-05 的基础上把站位横坐标作为连续决策变量，以 A 方案提供的可行覆盖包络作为约束，使用序列二次规划或信赖域方法减少局部漏测和过度重叠。第三，若获得船舶最小转弯半径、速度和加速度参数，可把曲率约束写成 $R(s)\ge R_{\mathrm{ship}}$，并把测线间连接航段纳入完整航次成本。第四，可对深度网格加入测量误差区间，以最不利深度和坡度重新评估覆盖，从而得到鲁棒测线方案。

\subsection{四问模型的最终归纳}
表~\ref{tab:final-summary} 用“模型—算法—结论—边界”四列概括全文，使评阅者能够快速识别每一问的数学核心及其适用范围。

\begin{table}[H]
\centering
\caption{四问模型、算法、主要结论与适用边界}
\label{tab:final-summary}
\begin{tabularx}{\textwidth}{C{1.1cm}C{3.2cm}C{3.2cm}L L}
\toprule
问题 & 数学模型 & 求解算法 & 主要结论 & 适用边界 \\
\midrule
一 & 斜坡声束交点与有向重叠率 & 显式逐线计算 & 浅水侧 $600,800$ m 处漏测 & 均匀平面坡面、固定航向 \\
二 & 深度梯度的航向/横航向分解 & $8\times8$ 参数枚举 & 得到完整覆盖宽度矩阵和方向对称性 & 均匀坡面、局部直线航向 \\
三 & 首线边界方程与最大间距显式递推 & 逐线递推至东边界 & $34$ 条，\SI{125936}{m}；$33$ 条不足 & 仅限等深线平行直线族 \\
四 & 双线性地形、局部条带、有限曲率曲线与离散指标 & 二分布线、站位插值、Bézier 平滑、独立评估 & 推荐 $61$ 条，\SI{576726.636}{m}，名义漏测 $0$ & 采样证据、有限候选族、无船型阈值 \\
\bottomrule
\end{tabularx}
\end{table}
